\begin{frame}{GitLab là gì? SaaS vs self-hosted}
\begin{itemize}
\item Nền tảng DevOps all-in-one trên Git: repo, Merge Request, Issues, CI/CD tích hợp sẵn.
\item Một vòng đời khép kín: plan, code, review, test, release, monitor trong cùng project.
\item Dùng ngay bản SaaS tại \texttt{gitlab.com} hoặc tự vận hành bản self-hosted.
\end{itemize}
\begin{center}
\small
\begin{tabular}{lll}
\toprule
Tiêu chí & SaaS & Self-hosted \\
\midrule
Hạ tầng & GitLab quản lý & Nhóm tự cài đặt \\
Chi phí \& bảo trì & Trả theo gói, không lo máy chủ & Miễn phí bản Community, tốn công admin \\
Phù hợp & Đồ án môn học, demo nhanh & Doanh nghiệp cần kiểm soát dữ liệu \\
\bottomrule
\end{tabular}
\end{center}
\end{frame}

\begin{frame}{Groups, Projects, Issues, Boards}
\begin{itemize}
\item Groups gom nhiều project theo nhóm/lớp: phân quyền member, chia sẻ runner chung.
\item Projects chứa repo Git, Merge Request, pipeline, wiki; gắn ASPX qua \texttt{main} bảo vệ.
\item Issues theo dõi task/bug với label, milestone, assignee; đóng bằng \texttt{Closes} khi merge.
\item Boards kiểu Kanban kéo thả Issues qua các cột To Do, Doing, Done theo milestone.
\end{itemize}
\end{frame}

\begin{frame}{Merge Request và Approval rules}
\begin{enumerate}
\item Tạo nhánh \texttt{feature} từ \texttt{main}, push và mở Merge Request kèm mô tả rõ ràng.
\item Hệ thống chạy pipeline kiểm tra build và test tự động trên nhánh nguồn.
\item Reviewer nhận xét từng dòng, yêu cầu chỉnh sửa cho đến khi đạt yêu cầu.
\item Đủ lượt approve theo rule mới được phép merge, đảm bảo ít nhất 2 người đồng ý.
\item Nhấn merge, xóa nhánh nguồn, Issues liên quan tự đóng theo từ khóa.
\end{enumerate}
\vspace{0.3em}
Ví dụ rule: nhánh \texttt{main} yêu cầu 2 approve từ maintainer và pipeline phải xanh 100\%.
\end{frame}

\begin{frame}[fragile]{Tệp gitlab-ci YAML cho ASPX 4.8.1}
\begin{itemize}
\item Khai báo \texttt{stages}, \texttt{variables}, job theo từng giai đoạn build, test, deploy.
\end{itemize}
\begin{lstlisting}[language=bash]
stages:
  - build
  - test
  - deploy
variables:
  SOLUTION: AspxDemo.sln
  CONFIG: Release
build:aspx:
  stage: build
  tags: [windows-481]
  script:
    - nuget restore %SOLUTION%
    - msbuild %SOLUTION% /p:Config=%CONFIG%
  artifacts:
    paths: [bin/Release/]
    expire_in: 1 week
test:aspx:
  stage: test
  tags: [windows-481]
  script: [vstest.console.exe bin/Release/*.Tests.dll]
\end{lstlisting}
\end{frame}

\begin{frame}{Runners và tags Windows}
\begin{itemize}
\item Runner là agent chạy job: shell runner trên Windows Server cài Visual Studio và \texttt{msbuild}.
\item Tags điều phối job đúng máy: job ASPX 4.8.1 gắn tag \texttt{windows-481}, runner cùng tag mới nhận.
\item Runner nhóm dùng chung cho 3 thành viên; bật cache NuGet để build lần sau nhanh hơn.
\end{itemize}
\end{frame}

\begin{frame}{Environments và Registry}
\begin{itemize}
\item Environments quản lý staging và production, hỗ trợ deploy thủ công và rollback một nút bấm.
\item Review Apps tự tạo môi trường xem trước cho mỗi Merge Request, xóa khi merge xong.
\item Container Registry lưu image Docker và gói NuGet nội bộ, dùng token nhóm để pull an toàn.
\end{itemize}
\end{frame}

\begin{frame}{Demo TV3: MR và pipeline}
\begin{itemize}
\item[$\square$] Mirror repo ASPX sang GitLab, bảo vệ nhánh \texttt{main} và bật Approval 2 người.
\item[$\square$] Mở Merge Request từ \texttt{feature} kèm mô tả, label và người review cụ thể.
\item[$\square$] Theo dõi pipeline build \texttt{msbuild} trên runner Windows đến khi báo xanh.
\item[$\square$] Thu đủ approve, merge vào \texttt{main} và kiểm tra deploy staging thành công.
\item[$\square$] Chụp pipeline xanh và bảng so sánh đưa vào báo cáo Chương 5 và Chương 6.
\end{itemize}
\end{frame}
